\section{Background and Review Method}
\label{sec:background}

\subsection{Signs of drowsiness}
An explanation can only be judged against what is already known about drowsiness. Four groups of signs recur in the literature \cite{fu2024survey,ali2026intelligent}.

\textit{Brain activity.} Theta and alpha power generally rise as alertness falls. Alpha spindles and theta bursts are the transient forms most often linked to drowsiness in the EEG studies reviewed here \cite{cui2022eeg,cui2022compact}.

\textit{Heart activity.} Heart rate tends to fall with low arousal, and heart rate variability (HRV) measures such as RMSSD, high-frequency power and the LF/HF ratio are read as indicators of autonomic balance \cite{puspasari2026predicting}.

\textit{Eyes, face and head.} Longer and slower blinks, a higher percentage of eye closure over time (PERCLOS), yawning and head nodding are the standard behavioural signs \cite{fu2024survey,hasan2024validation}.

\textit{Vehicle control.} Fewer small steering corrections, sometimes followed by a large correction, can precede a drowsiness-related lane departure \cite{mcdonald2014steering}.

\subsection{Interpretable versus explainable}
Many papers use the two terms interchangeably \cite{xaieeg2026review}. We separate them by mechanism. A method is \textit{intrinsic} (interpretable by design) when the decision can be read directly from the model: a rule, a threshold, a transparent feature or a model's own internal state. A method is \textit{post-hoc} (explainable) when a separate computation is run after training to attribute the decision to the input. This follows Rudin's distinction between explaining a black box and using an interpretable model \cite{rudin2019stop}. It also means we sometimes classify a paper differently from how its authors describe it. For example, we treat class activation mapping as post-hoc and model-specific, even when the network was designed to support it, because the map is computed after training.

\subsection{Review method}
This is a structured narrative review, not a systematic review under PRISMA. We searched with general academic web search, Europe PMC and arXiv, and checked IEEE Xplore for overlapping reviews. The base query combined three groups of terms: (drowsiness, drowsy, sleepiness, fatigue, vigilance); (driver, driving, driver monitoring); and (explainable, interpretable, XAI, SHAP, LIME, Grad-CAM, class activation, saliency, layer-wise relevance, integrated gradients, attention). Searches were repeated with added terms for each modality group.

We selected 30 papers, ten per modality group, that use an interpretable model or apply an explanation method to drowsiness or fatigue detection, preferring recent work (2024--2026) and including older papers where they are foundational. We recorded no screening counts, so we cannot report how many records were identified or excluded.

Each paper was read in full by one author, who completed a fixed 12-field template covering data, labels, model, explanation method, findings, results, limitations and four appraisal checks. Key numbers and the strongest critical statements were then checked against the source PDFs. Three papers are excluded from the appraisal counts: a methods study \cite{cui2023towards}, a dataset paper with no explanation method \cite{ghoddoosian2019realistic}, and a study that classifies driving manoeuvres rather than drowsiness \cite{askari2026physiologically}. That leaves 27 studies, 19 of them from 2024--2026.

\textit{Limits of this method.} The selection is targeted, not exhaustive, and it includes weak studies deliberately so that we could examine the range of practice. The counts therefore describe the studies we examined and should not be read as prevalence across the field. Each paper had one reader, and the appraisal involves judgement; the per-study ratings in Tables~\ref{tab:eeg}--\ref{tab:vision} are given so that readers can check them.
